\def\confversion{0}
\def\ifconf{\ifnum\confversion=1}
\def\ifnotconf{\ifnum\confversion=0}
\def\showauthornotes{0}
\def\showkeys{0}
\def\showdraftbox{0}
%%%%%%%%%%%% Packages
\usepackage{xspace,xcolor,enumerate}
\usepackage{amsmath,amssymb}
\usepackage{amsthm}
\usepackage[toc,page]{appendix}
\usepackage{thmtools}
\usepackage{thm-restate}
\usepackage{color,graphicx}
\usepackage{makecell}
\usepackage{tabularx}
\usepackage{enumitem}
%\usepackage{prettyref}
% \usepackage[varg]{txfonts} % varg - uses nicer g,v,y,w
\usepackage[linesnumbered,ruled,vlined]{algorithm2e}
\ifnum\showkeys=1
\usepackage[color]{showkeys}
\fi

\definecolor{darkred}{rgb}{0.5,0,0}
\definecolor{darkgreen}{rgb}{0,0.35,0}
\definecolor{darkblue}{rgb}{0,0,0.55}

\usepackage[pdfstartview=FitH,pdfpagemode=UseNone,colorlinks,linkcolor=darkblue,filecolor=darkred,citecolor=darkgreen,urlcolor=darkred,pagebackref]{hyperref}

\usepackage[capitalise,nameinlink]{cleveref}
\usepackage[T1]{fontenc}
\usepackage{mathtools,dsfont,bbm}
%\usepackage{mathpazo} % Pei: this font is very nice on text, but does not feel right to me for math formulas
%\usepackage[top=1in, bottom=1in, left=1in, right=1in]{geometry}
% \usepackage[top=1in, bottom=1in, left=1.25in, right=1.25in]{geometry}
% \usepackage{fullpage}


%%%%%%%%%%%%%%% Lengths
% \setlength{\parindent}{0 in} 
\setlength{\parskip}{0.05 in}
\setlength{\parindent}{4 ex} 
% \setlength{\parskip}{0.5 ex}

%%%%%%%%%%%%%%% Author Notes
\ifnum\showauthornotes=1
\newcommand{\Authornote}[2]{{\sf\small\color{red}{[#1: #2]}}}
\newcommand{\Authorcomment}[2]{{\sf \small\color{gray}{[#1: #2]}}}
\newcommand{\Authorfnote}[2]{\footnote{\color{red}{#1: #2}}}
\else
\newcommand{\Authornote}[2]{}
\newcommand{\Authorcomment}[2]{}
\newcommand{\Authorfnote}[2]{}
\fi

%%%%%%%%%%%%%%%% Draftbox
\ifnum\showdraftbox=1
\newcommand{\draftbox}{\begin{center}
  \fbox{%
    \begin{minipage}{2in}%
      \begin{center}%
        \begin{Large}%
          \textsc{Working Draft}%
        \end{Large}\\
        Please do not distribute%
      \end{center}%
    \end{minipage}%
  }%
\end{center}
\vspace{0.2cm}}
\else
\newcommand{\draftbox}{}
\fi

%%%%%%%%%%%%%% No bullets
\renewcommand{\labelitemi}{-}


%%%%%%%%%%%%%%%%%% Theorem Environments
\newtheorem{theorem}{Theorem}[section]
\newtheorem{conjecture}[theorem]{Conjecture}
\newtheorem{hypothesis}[theorem]{Hypothesis}
\newtheorem{observation}[theorem]{Observation}
\newtheorem{definition}[theorem]{Definition}
\newtheorem{notation}[theorem]{Notation}
\newtheorem{openproblem}[theorem]{Open Problem}
\newtheorem{lemma}[theorem]{Lemma}
\newtheorem{remark}[theorem]{Remark}
\newtheorem{proposition}[theorem]{Proposition}
\newtheorem{corollary}[theorem]{Corollary}
\newtheorem{claim}[theorem]{Claim}
\newtheorem{fact}[theorem]{Fact}
\newtheorem{openprob}[theorem]{Open Problem}
\newtheorem{remk}[theorem]{Remark}
\newtheorem{example}[theorem]{Example}
\newtheorem{apdxlemma}{Lemma}
%\newtheorem{algorithm_simple}[theorem]{Algorithm}
%[Tushant]Changed algo env to remark style - no italics - revert later if needed
\theoremstyle{remark}
\newtheorem{algo}[theorem]{Algorithm}
\newcommand{\question}[1]{{\sf [#1]\marginpar{?}} }

%Commented this out - Tushant

\newenvironment{algorithm_simple}[3]
        {\noindent\begin{boxedminipage}{\textwidth}
        %\theorembodyfont{\rm}
        \begin{algo}[#1]\ \par
        {\begin{tabular}{r l}
        \textbf{Input} & #2\\
        \textbf{Output} & #3
        \end{tabular}\par\enskip}}
        {\end{algo}\end{boxedminipage}}

%%%%%%%%%%%%%%%%% Proof Environments

\def\FullBox{\hbox{\vrule width 6pt height 6pt depth 0pt}}

\iffalse
\def\qed{\ifmmode\qquad\FullBox\else{\unskip\nobreak\hfil
\penalty50\hskip1em\null\nobreak\hfil\FullBox
\parfillskip=0pt\finalhyphendemerits=0\endgraf}\fi}

\fi
\def\qedsketch{\ifmmode\Box\else{\unskip\nobreak\hfil
\penalty50\hskip1em\null\nobreak\hfil$\Box$
\parfillskip=0pt\finalhyphendemerits=0\endgraf}\fi}
%\ifnotconf
%\renewenvironment{proof}{\begin{trivlist} \item {\bf Proof:~~}}
%   {\qed\end{trivlist}}
%\fi

\newenvironment{proofsketch}{\begin{trivlist} \item {\bf
Proof Sketch:~~}}
  {\qedsketch\end{trivlist}}

\newenvironment{proofof}[1]{\begin{trivlist} \item {\bf Proof of
#1:~~}}
  {\qed\end{trivlist}}

\newenvironment{claimproof}{\begin{quotation} \noindent
{\bf Proof of claim:~~}}{\qedsketch\end{quotation}}


%%%%%%%%%% Symbols and Fonts
\newcommand{\acts}{{\cdot}}
\def\from{:}
\def\to{\rightarrow}
\def\eps{\varepsilon}
\def\epsilon{\varepsilon}
\def\e{\epsilon}
\def\eps{\epsilon}
\def\d{\delta}
%\def\phi{\varphi} %Pei varphi and psi are two similar
\def\cal{\mathcal}
\def\xor{\oplus}
\def\ra{\rightarrow}
\def\implies{\Rightarrow}
\newcommand{\sub}{\ensuremath{\subseteq}}
\def\psdgeq{\succeq} 
% \newcommand{\defeq}{\stackrel{\mathrm{def}}=}     
\newcommand{\defeq}{:=}     
\providecommand{\cod}{\mathop{\sf cod}\nolimits}


\renewcommand{\bar}{\overline} 
\newcommand{\ol}{\overline}
% \renewcommand{\mathbb}{\varmathbb}
% \renewcommand{\Bbbk}{\varBbbk}

\def\bull{\vrule height .9ex width .8ex depth -.1ex }

%%%%%%%%%%%%%%%%%%%%%% Text Macros
\newcommand{\ie}{i.e.,\xspace}
\newcommand{\eg}{e.g.,\xspace}
\newcommand{\etal}{et al.\xspace}
\newcommand{\cf}{{\it cf.,}}

%%%%%%%%%%%%%%%%%%%%% Punctuation at the end of a displayed formula
\newcommand{\mper}{\,.}
\newcommand{\mcom}{\,,}

%%%%%%%%%%%%%%%%%%%%%% Number Sets
\newcommand{\R}{{\mathbb R}}
\newcommand{\E}{{\mathbb E}}
\newcommand{\C}{{\mathbb C}}
\newcommand{\N}{{\mathbb{N}}}
\newcommand{\Z}{{\mathbb Z}}
\newcommand{\Q}{{\mathbb{Q}}}
\newcommand{\F}{{\mathbb F}}
\newcommand{\zo}{\{0,1\}}
\newcommand{\GF}{\mathrm{GF}}
\newcommand{\FF}{{\mathbb F}}
\newcommand{\Real}{{\mathbb R}}
\newcommand{\cA}{\mathcal{A}}
\newcommand{\cB}{\mathcal{B}}
\newcommand{\cC}{\mathcal{C}}
\newcommand{\cH}{\mathcal{H}}
\newcommand{\cD}{\mathcal{D}}
\newcommand{\cF}{\mathcal{F}}
\newcommand{\cE}{\mathcal{E}}
\newcommand{\cN}{\mathcal{N}}
\newcommand{\cG}{\mathcal{G}}
\newcommand{\cU}{\mathcal{U}}
\newcommand{\cS}{\mathcal{S}}
\newcommand{\cT}{\mathcal{T}}


\newcommand{\B}{\{0,1\}\xspace}
\newcommand{\pmone}{\{-1,1\}\xspace}
\newcommand{\on}{\pmone}
\newcommand{\hcube}{\ensuremath{\{-1,1\}^N}}
\newcommand{\Zk}{\mathbb{Z}_k}
\newcommand{\indicator}[1]{\mathds{1}_{\{#1\}}}

%%%%%%%%%% Standard Normal Distribution

\newcommand{\gauss}[2]{{\cal N(#1, #2)}}
\newcommand{\gaussian}[2]{{\cal N(#1, #2)}}

%%%%%%%%%%%% Fractions
% commands for fractions 
\usepackage{nicefrac}
% poor man's fraction
\newcommand{\flatfrac}[2]{#1/#2}
\newcommand{\varflatfrac}[2]{#1\textfractionsolidus#2}

\let\nfrac=\nicefrac
\let\ffrac=\flatfrac
% similar commands: tfrac,dfrac

\newcommand{\half}{\nicefrac12}
\newcommand{\onequarter}{\nicefrac14}
\newcommand{\threequarter}{\nicefrac34}

%%%%%%%%%%%%%%%%%% Vectors and Reals
\newcommand{\abs}[1]{\ensuremath{\left\lvert #1 \right\rvert}}
\newcommand{\smallabs}[1]{\ensuremath{\lvert #1 \rvert}}
%
%------------------------------------ Norms -------------------
\newcommand{\norm}[1]{\ensuremath{\left\lVert #1 \right\rVert}}
\newcommand{\smallnorm}[1]{\ensuremath{\lVert #1 \rVert}}


% \makeatletter
% \def\norm#1{%fancy norm command
%   \@ifnextchar\bgroup%
%    {\normalpnorm{#1}}%
%    {\defaultnorm{#1}}%
% }
% \def\defaultnorm#1{%
%     \renderNorm{#1}{}
% }
% \def\normalpnorm#1#2{%
%    \@ifnextchar\bgroup%
%    {\banaxnorm{#1}{#2}}%
%    {\renderNorm{#2}{#1}}%\ensuremath{\left\lVert {#2} \right\rVert_{{#1}}}}%
% }

% \def\banaxnorm#1#2#3{%
%     \renderNorm{#3}{#1\to#2 }
%     %\ensuremath{\left\lVert #3 \right\rVert_{{#1}\to{#2}}}%
% }

% \def\renderNorm#1#2{%
%     \@ifnextchar^%
%     {\fixExponent{#1}{#2}}%
%     {\ensuremath{\mathchoice%
%         {\lVert #1 \rVert_{#2}}%
%         {\lVert #1 \rVert_{#2}}%
%         {\lVert #1 \rVert_{#2}}%
%         {\lVert #1 \rVert_{#2}}}%
%     }%
% }
% \def\fixExponent#1#2^#3{%
%     \ensuremath{{\lVert #1 \rVert^{#3}_{#2}}}%
% }

% \def\enorm#1#2{%
%    \@ifnextchar\bgroup%
%    {\norm{L_{#1}}{L_{#2}}}%
%    {\norm{L_{#1}}{#2}}%
% }

% \def\cnorm#1#2{%
%    \@ifnextchar\bgroup%
%    {\norm{\ell_{#1}}{\ell_{#2}}}%
%    {\norm{\ell_{#1}}{#2}}%
% }

% \makeatother

%
\newcommand{\mydot}[2]{\ensuremath{\left\langle #1, #2 \right\rangle}}
\newcommand{\mysmalldot}[2]{\ensuremath{\langle #1, #2 \rangle}}
%\newcommand{\ip}[1]{\left\langle #1 \right\rangle}
\newcommand{\ip}[2] {\ensuremath{\left\langle #1 , #2 \right\rangle}}

\newcommand{\emysmalldot}[2]{\ensuremath{\langle #1, #2 \rangle}_E} %% Euiwoong 20171030
\newcommand{\cmysmalldot}[2]{\ensuremath{\langle #1, #2 \rangle}_C} %% Euiwoong 20171030

\newcommand{\cip}[1]{\left\langle #1 \right\rangle}  %% Euiwoong 20171030
\newcommand{\eip}[1]{\left\langle #1 \right\rangle} %% Euiwoong 20171030

%%%%%%%%%%% Vectors
\def\bfx {{\bf x}}
\def\bfy {{\bf y}}
\def\bfz {{\bf z}}
%
\def\bfu{{\bf u}}
\def\bfv{{\bf v}}
\def\bfw{{\bf w}}
%
\def\bfa{{\bf a}}
\def\bfb{{\bf b}}
\def\bfc{{\bf c}}
%
\def\bfg{{\bf g}}

\newcommand{\zero}{\mathbf 0}
\newcommand{\one}{{\mathbf{1}}}
\newcommand{\zeroone}{{0/1}\xspace}
\newcommand{\minusoneone}{{-1/1}\xspace}

\newcommand{\yes}{{\sf Yes}\xspace}
\newcommand{\no}{{\sf No}\xspace}

%%%%%%%%%%%%%%%%%%%%%%%% Lasserre Variables
\newcommand{\vone}[1]{\mathbf{u}_{#1}}
\newcommand{\voneempty}{\mathbf{u}_{\emptyset}}
\newcommand{\vtwo}[2]{\mathbf{v}_{(#1, #2)}}
\newcommand{\vtwoempty}{\mathbf{v}_{(\emptyset, \emptyset)}}

%%%%%%%%%%%%%%%% Sherali-Adams
\newcommand{\qary}{[q]}
\newcommand{\varone}[1]{y_{#1}}
\newcommand{\varoneempty}{y_{\emptyset}}
\newcommand{\vartwo}[2]{x_{(#1, #2)}}
\newcommand{\vartwoempty}{x_{(\emptyset,\emptyset)}}


%%%%%%%%%%%%%%%%%%%%% Random Variables
\newcommand{\Esymb}{\mathbb{E}}
\newcommand{\Psymb}{\mathbb{P}}
\newcommand{\Vsymb}{\mathbb{V}}
\newcommand{\Varsymb}{\mathrm{Var}}
\DeclareMathOperator*{\ExpOp}{\Esymb}
\DeclareMathOperator*{\ProbOp}{\Psymb}

%%%% Variable arguments \Pr, \Ex, \Var, \tildeEx
%%%%  changes its behavior according to number of arguments as follows:
%%%%    - 1: <symbol> [ #1 ]
%%%%    - 2: <symbol>_[#1] [ #2 ] -- this behaves like \max, i.e., subscript is placed
%%%%                underneath the symbol if in displaystyle environment
%%%%    - 4: <symbol>_[#1]  #3 #2 #4 -- here #3 and #4 are the opening and closing
%%%%                braces symbol
\makeatletter
%\def\Pr#1{%
%    \ProbabilityRender{\Psymb}{#1}%
%}

\def\Ex#1{%
    \ProbabilityRender{\Esymb}{#1}%
}

\def\tildeEx#1{%
    \ProbabilityRender{\widetilde{\Esymb}}{#1}%
}

\def\Var#1{%
    \ProbabilityRender{\Varsymb}{#1}%
}

\def\ProbabilityRender#1#2{%fancy probability command
  \@ifnextchar\bgroup%
  {\renderwithdist{#1}{#2}}
   {\singlervrender{#1}{#2}}
}
\def\singlervrender#1#2{%
   \ensuremath{\mathchoice
       {{#1}\left[ #2 \right]}
       {{#1}[ #2 ]}
       {{#1}[ #2 ]}
       {{#1}[ #2 ]}
   }
}
\def\renderwithdist#1#2#3{%
   \@ifnextchar\bgroup
   {\superfancyrender{#1}{#2}{#3}}
   {\ensuremath{\mathchoice
      {\underset{#2}{#1}\left[ #3 \right]}
      {{#1}_{#2}[ #3 ]}
      {{#1}_{#2}[ #3 ]}
      {{#1}_{#2}[ #3 ]}
     }
   }
}
\def\superfancyrender#1#2#3#4#5{
   \ensuremath{\mathchoice
      {\underset{#1}{{#1}}\left#4 #3 \right#5}
      {{#1}_{#2}#4 #3 #5}
      {{#1}_{#2}#4 #3 #5}
      {{#1}_{#2}#4 #3 #5}
   }
}
\makeatother


\def\expop{\ExpOp}
\def\probop{\ProbOp}

\newcommand\restrict[1]{\raisebox{-.5ex}{$|$}_{#1}}

\newcommand{\conv}[1]{\mathrm{conv}\inparen{#1}}


%\newcommand{\varex}[1]{\E\paren{#1}}
%\newcommand{\varEx}[1]{\E\Paren{#1}}

\newcommand{\given}{\;\middle\vert\;}


%%%%%%%%%%%%%%%%%%%%%% General Useful Macros
\newfont{\inhead}{eufm10 scaled\magstep1}
\newcommand{\deffont}{\sf}

\newcommand{\calA}{{\cal A}}
\newcommand{\calB}{{\cal B}}
\newcommand{\calC}{{\cal C}}
\newcommand{\calD}{{\cal D}}
\newcommand{\calE}{{\cal E}}
\newcommand{\calF}{{\cal F}}
\newcommand{\calG}{{\cal G}}
\newcommand{\calH}{{\cal H}}
\newcommand{\calL}{{\cal L}}
\newcommand{\calU}{{\cal U}}
\newcommand{\calP}{{\cal P}}
\newcommand{\calO}{{\cal O}}
\newcommand{\calX}{{\cal X}}
\newcommand{\calS}{{\cal S}}

% \newcommand{\maximize}{\mathop{\textrm{maximize}}}
% \newcommand{\minimize}{\mathop{\textrm{minimize}}}
% \newcommand{\subjectto}{\mathop{\textrm{subject to}}}


\newcommand{\poly}{{\mathrm{poly}}}
\newcommand{\polylog}{{\mathrm{polylog}}}
\newcommand{\loglog}{{\mathop{\mathrm{loglog}}}}
\newcommand{\suchthat}{{\;\; : \;\;}}
\newcommand{\getsr}{\mathbin{\stackrel{\mbox{\tiny R}}{\gets}}}


%%%%%%%%%% Operators
\DeclareMathOperator\supp{supp}

%%%%%%%%%% Asymptotic and math symbols borrowed from Laci
\DeclareMathOperator{\Div}{\operatorname {Div}}
\DeclareMathOperator{\cov}{\operatorname {Cov}}
\DeclareMathOperator{\period}{\operatorname {per}}
\DeclareMathOperator{\rank}{\operatorname {rk}}
\DeclareMathOperator{\tr}{\operatorname {tr}}
\DeclareMathOperator{\Span}{\operatorname {span}}
\DeclareMathOperator{\sgn}{\operatorname{sgn}}
\DeclareMathOperator{\inv}{\operatorname{Inv}}
\DeclareMathOperator{\sym}{\operatorname{Sym}}
\DeclareMathOperator{\alt}{\operatorname{Alt}}
\DeclareMathOperator{\aut}{\operatorname{Aut}}
\DeclareMathOperator{\ord}{\operatorname{ord}}
\DeclareMathOperator{\lcm}{\operatorname{lcm}}
\DeclareMathOperator*{\argmin}{\arg\!\min}
\DeclareMathOperator*{\argmax}{\arg\!\max}
\DeclareMathOperator{\Paley}{\operatorname {Paley}}
\DeclareMathOperator{\Vtx}{\operatorname V}
\DeclareMathOperator{\Edg}{\operatorname E}
\DeclareMathOperator{\im}{\operatorname{im}}
\DeclareMathOperator{\ran}{\operatorname{ran}}
\DeclareMathOperator{\diag}{\operatorname{diag}}
\DeclareMathOperator{\erf}{\operatorname{erf}}
\newcommand{\ee}{\ensuremath{\mathrm e}}
\newcommand{\w}{\ensuremath{\omega}}
\newcommand{\vp}{\ensuremath{\varphi}}
\newcommand{\ve}{\ensuremath{\varepsilon}}
\newcommand{\el}{\ensuremath{\ell} }
\newcommand{\CCC}{\ensuremath{\mathbb{C}}}
\newcommand{\EEE}{\ensuremath{\mathbb{E}}}
\newcommand{\FFF}{\ensuremath{\mathbb{F}}}
\newcommand{\QQQ}{\ensuremath{\mathbb{Q}}}
\newcommand{\RRR}{\ensuremath{\mathbb{R}}}
\newcommand{\NNN}{\ensuremath{\mathbb{N}}}
\newcommand{\ZZZ}{\ensuremath{\mathbb{Z}}}
%\renewcommand{\iff}{\ensuremath{\Leftrightarrow}}
\renewcommand{\bar}[1]{\ensuremath{\overline{#1}}}
\newcommand{\ceil}[1]{\ensuremath{\left\lceil #1 \right\rceil}}
\newcommand{\floor}[1]{\ensuremath{\left\lfloor #1 \right\rfloor}}
%\newcommand{\1}[1]{\mathds{1}\left[#1\right]}
\newcommand{\1}[1]{\mathds{1}{[#1]}}

\newcommand{\bigoh}{\operatorname{O}}
\newcommand{\bigtheta}{\mathop{\Theta}}
\newcommand{\bigomega}{\mathop{\Omega}}
\newcommand{\littleoh}{\operatorname{o}}

%%%%%%%%%%%%%%%%%%% Complexity Classes
\newcommand{\classfont}[1]{\textsf{#1}}
\newcommand{\coclass}[1]{\mathbf{co\mbox{-}#1}} 
\newcommand{\BPP}{\classfont{BPP}\xspace}
\newcommand{\classP}{\classfont{P}\xspace}
%\newcommand{\NP}{\classfont{NP}\xspace}
\newcommand{\coNP}{\classfont{coNP}\xspace}
\newcommand{\nphard}{\classfont{NP}-hard\xspace}

%%%%%%%%%%%%%%%%%% Parameterized Complexity Classes 
\newcommand{\DTIME}[1]{\classfont{DTIME}\inparen{#1}\xspace}
\newcommand{\BPTIME}[1]{\classfont{BPTIME}\inparen{#1}\xspace}

%%%%%%%%%%%%%% SDPs and LPs
\newcommand{\opt}{{\sf OPT}\xspace}
\newcommand{\sdpopt}{{\sf FRAC}\xspace}
\newcommand{\lpopt}{{\sf FRAC}\xspace}
\newcommand{\round}{{\sf ROUND}\xspace}
% \newcommand{\fractional}{{\sf FRAC}}
% \newcommand{\minopt}{{\sf MINOPT}}
% \newcommand{\greedy}{{\sf Greedy}}



\newcommand{\obj}{\mbox{\sf obj}}
\newcommand{\sdp}{\mbox{\sf sdp}}
\newcommand{\val}{{\sf val}}

% \newcommand{\SDP}{{\sf SDP}\xspace}
% \newcommand{\LP}{{\sf LP}\xspace}

%%%%%%%%%%%%%%%%%%%% Problems
% use texorpdfstring to avoid problems with hyperref (can use problem
% macros also in headings
\newcommand{\problemmacro}[1]{\textsf{#1}}

\newcommand{\maxcsp}{\problemmacro{MAX-CSP}\xspace}
\newcommand{\maxkcsp}{\problemmacro{MAX k-CSP}\xspace}
\newcommand{\maxkcspq}{\problemmacro{MAX k-CSP}$_q$\xspace}
\newcommand{\maxksat}{\problemmacro{MAX k-SAT}\xspace}
\newcommand{\maxthreesat}{\problemmacro{MAX 3-SAT}\xspace}
\newcommand{\maxthreexor}{\problemmacro{MAX 3-XOR}\xspace}
\newcommand{\maxtwoxor}{\problemmacro{MAX 2-XOR}\xspace}
\newcommand{\maxfourxor}{\problemmacro{MAX 4-XOR}\xspace}
\newcommand{\maxkxor}{\problemmacro{MAX k-XOR}\xspace}
\newcommand{\maxcut}{\problemmacro{MAX-CUT}\xspace}

\newcommand{\independentset}{\problemmacro{Maximum Independent Set}\xspace}
\newcommand{\stableset}{\problemmacro{Maximum Stable Set}\xspace}
\newcommand{\vertexcover}{\problemmacro{Minimum Vertex Cover}\xspace}
\newcommand{\clique}{\problemmacro{Maximum Clique}\xspace}
\newcommand{\apxcoloring}{\problemmacro{Approximate Graph Coloring}\xspace}
\newcommand{\chromaticnumber}{\problemmacro{Chromatic Number}\xspace}
\newcommand{\threecolor}{\problemmacro{3-Coloring}\xspace}
\newcommand{\res}{{\sf Res}}

\newcommand{\sa}{\textrm{Sherali-Adams}}
\newcommand{\ls}{\textrm{Lov\'asz-Schrijver}}
\newcommand{\la}{\textrm{Lasserre}}
\newcommand{\gc}{\textrm{Gomory-Chv\'atal}}

\newcommand{\sparsestcut}{\problemmacro{Sparsest Cut}\xspace}
\newcommand{\conductance}{\problemmacro{Graph Conductance}\xspace}
\newcommand{\balancedseparator}{\problemmacro{Balanced Separator}\xspace}
\newcommand{\labelcover}{\problemmacro{Label Cover}\xspace}                                %% Added by Euiwoong on 2017-10-21.
\newcommand{\uniquegames}{\problemmacro{Unique Games}\xspace}
\newcommand{\ugc}{\problemmacro{Unique Games Conjecture}\xspace}
\newcommand{\mla}{\problemmacro{Minimum Linear Arrangement}\xspace}
\newcommand{\gammamaxlin}{\Gamma \problemmacro{MAX2LIN}\xspace}
\newcommand{\tsp}{\problemmacro{Traveling Salesman Problem}\xspace}
\newcommand{\maximummatching}{\problemmacro{Maximum Matching}\xspace}
\newcommand{\maxminallocation}{\problemmacro{MaxMin Allocation}\xspace}
\newcommand{\hypergraphindset}{\problemmacro{Hypergraph Independent Set}\xspace}
\newcommand{\knapsack}{\problemmacro{Knapsack}\xspace}

% Euiwoong Definitions
\newcommand{\threesat}{3-SAT\xspace}
\newcommand{\groth}[2]{\ensuremath{(#1, #2)}-Grothendieck}
\newcommand{\onorm}[2]{\ensuremath{#1{\to}#2} norm}



%Temporary [Tushant]

\renewcommand{\Re}{\operatorname{Re}}
\renewcommand{\Im}{\operatorname{Im}}



%%%%%%%%%%%%%%%%%%%%%%%%% Enclosures
\newcommand{\inparen}[1]{\left(#1\right)}             %\inparen{x+y}  is (x+y)
\newcommand{\inbraces}[1]{\left\{#1\right\}}           %\inbrace{x+y}  is {x+y}
\newcommand{\insquare}[1]{\left[#1\right]}             %\insquare{x+y}  is [x+y]
\newcommand{\inangle}[1]{\left\langle#1\right\rangle} %\inangle{A}    is <A>


%%%%%%%%%%%%%%%%%%%% Names
\newcommand{\Erdos}{Erd\H{o}s\xspace}
\newcommand{\Renyi}{R\'enyi\xspace}
\newcommand{\Lovasz}{Lov\'asz\xspace}
\newcommand{\Juhasz}{Juh\'asz\xspace}
\newcommand{\Bollobas}{Bollob\'as\xspace}
\newcommand{\Furedi}{F\"uredi\xspace}
\newcommand{\Komlos}{Koml\'os\xspace}
\newcommand{\Luczak}{\L uczak\xspace}
\newcommand{\Kucera}{Ku\v{c}era\xspace}
\newcommand{\Szemeredi}{Szemer\'edi\xspace}
\newcommand{\Chvatal}{Chv\'atal}
\newcommand{\Holder}{H\"{o}lder}
\newcommand{\Plunnecke}{Pl\"unnecke}
\newcommand{\Odonnell}{O'Donnell}
\newcommand{\Briet}{Bri\"et\xspace}          %% Added by Euiwoong on 2017-10-22.
\newcommand{\Esseen}{Ess\'een\xspace}          %% Added by Euiwoong on 2017-10-22.



%%%%%%%%%%%%%%%%%%%%%% Quantum
\newcommand{\ket}[1]{\lvert #1\rangle}
\newcommand{\braket}[2]{\left\langle #1 \,\middle\vert\, #2\right\rangle}
\newcommand{\QMA}{\textup{QMA}}
\newcommand{\BQP}{\textup{BQP}}
\newcommand{\NP}{\textup{NP}}
\newcommand{\NEXP}{\textup{NEXP}}
\newcommand{\PCP}{\textup{PCP}}
\newcommand{\Pclass}{\textup{P}}

% filled card suits
\DeclareSymbolFont{extraup}{U}{zavm}{m}{n}
\DeclareMathSymbol{\varheart}{\mathalpha}{extraup}{86}
\DeclareMathSymbol{\vardiamond}{\mathalpha}{extraup}{87}

\def\RR{\mathbb R}
\def\totw{\text{ otherwise }}
\def\tif{\text{ if }}
\def\wio{\backslash}
\def\tand{\quad\text{and}\quad}
\def\tor{\text{ or }}
\def\tst{\text{ s.t. }}
\def\matr#1{\mathsf{#1}}
\def\vecc#1{\boldsymbol{#1}}
\def\rand#1{\mathbf{#1}}
\def\rv#1{\rand #1}
\def\conj{\overline}
\def\SymG{\mathfrak S}
\def\Sep{\mathcal S}
\def\Lap{\mathcal L}
\def\Ray{\mathcal R}
\def\Adj{\mathcal A}
\def\Ker{\mathcal K}
\def\Wal{\mathcal W}
\def\One{\mathbb 1}
\def\one{\mathbf 1}
\def\Sp{\mathbb S}
\def\Ba{\mathbb B}
\def\Aa{\mathcal A}
\def\Pred{\mathcal P}
\def\Cc{\mathcal C}
\def\Xx{\mathcal X}
\def\Tt{\mathcal T}
\def\Ff{\mathcal F}
\def\Ee{\mathcal E}
\def\Di{\mathcal D}
\def\tensor{\otimes}
\def\ALG{\mathsf{ALG}}
\def\OPT{\mathsf{OPT}}
\def\SAT{\mathsf{SAT}}
\def\AVG{\mathsf{AVG}}
\def\SDP{\mathsf{SDP}}
\def\LP{\mathsf{LP}}
\def\vph{\varphi}
\def\ee{\varepsilon}
\def\pmu{\widetilde \mu}
\def\dif{\textrm{d}}
\def\Ins{\mathfrak I}
\def\assn{\sigma}
\def\tree{\mathcal T}
%% matrices
\def\Up{\matr U}
\def\Dee{\matr D}
\def\Aye{\matr A}
\def\Bee{\matr B}
\def\Why{\matr Y}
\def\Www{\matr W}
\def\Enn{\matr N}
\def\Ess{\matr S}
\def\Ide{\matr I}
\def\Pii{\matr P}
\def\Emm{\matr M}
\def\Jay{\matr J}
\def\Gee{\matr G}
%% functions
\def\eff{\vecc f}
\def\gee{\vecc g}
\def\eee{\vecc e}
\def\etch{\vecc h}
\def\pii{\vecc p}
\def\cee{\vecc c}
\def\wii{\vecc w}
%% faces
\def\aye{\mathfrak a}
\def\ess{\mathfrak s}
\def\tee{\mathfrak t}
\def\bee{\mathfrak b}
\def\teen{\widetilde{\tee}}
% operators
\DeclareMathOperator{\SwapT}{Swap}
\DeclareMathOperator{\Ent}{Ent}
\DeclareMathOperator{\hamw}{wt}
\DeclareMathOperator{\ddeg}{deg}
\DeclareMathOperator{\Det}{det}
\DeclareMathOperator{\sspan}{span}
\DeclareMathOperator{\Sym}{Sym}
\DeclareMathOperator{\Deg}{deg}
\DeclareMathOperator{\Herm}{\mathcal H}
\DeclareMathOperator{\Pos}{\mathcal P}
\DeclareMathOperator{\Dens}{\mathcal D}
\DeclareMathOperator{\OO}{\mathcal O}
\DeclareMathOperator{\OOs}{\widetilde{\mathcal O}}
\DeclareMathOperator{\Oms}{\widetilde{\Omega}}
\DeclareMathOperator{\oo}{o}
\DeclareMathOperator{\Tr}{Tr}
\DeclareMathOperator{\Cay}{Cay}
\DeclareMathOperator{\colp}{cp}
\DeclareMathOperator{\plog}{polylog}
\DeclareMathOperator{\vol}{vol}
\DeclareMathOperator{\ddiag}{diag}
\DeclareMathOperator{\MI}{I}
\DeclareMathOperator{\Stab}{Stab}
\DeclareMathOperator{\NoSe}{NS}
\DeclareMathOperator{\dist}{d}
\DeclareMathOperator{\parikh}{\#}
\DeclareMathOperator{\ddiam}{diam}
\DeclareMathOperator{\Num}{N}
\DeclareMathOperator{\spec}{Spec}
\DeclareMathOperator{\Diff}{\Delta}
\DeclareMathOperator{\Rank}{rank}
\DeclareMathOperator{\PExp}{\widetilde \E}
% parens, braces, brackets, norms, absolute value, and inner products
\DeclarePairedDelimiter\set{\lbrace}{\rbrace}
\DeclarePairedDelimiter\parens{\lparen}{\rparen}
\DeclarePairedDelimiter\sqbr{[ }{]}
\DeclarePairedDelimiter\Abs{|}{|}


\DeclareMathOperator{\bias}{bias}
\DeclareMathOperator{\lift}{dsum}
\DeclareMathOperator{\dsum}{dsum}
\DeclareMathOperator{\dprod}{dprod}

%%%%%%%%%%%%%%%%%%%%%%%%%%%%%%%%
%
%       Walk Operators
%
%%%%%%%%%%%%%%%%%%%%%%%%%%%%%%%%


%%%%%%%%%%%%%%%%%%%%%%%%%%%%%%%%
%        Swap Walks
%%%%%%%%%%%%%%%%%%%%%%%%%%%%%%%%

% Height Specific Rectangular Swap Walk
% Args:
%       #1='k' size of initial face s
%       #2='l' size of final face t
%       #3='j' number of new elements, i.e., |t\s|
%       #4='u' walk height
\newcommand{\hrswap}[4]{\Ess_{#1,#2}^{\left(#3,#4\right)}}

% Rectangular Swap Walk (Height Independent)
% Args:
%       #1='k' size of initial face s
%       #2='l' size of final face t
%       #3='j' number of new elements, i.e., |t\s|
\newcommand{\rswap}[3]{\Ess_{#1,#2}^{\left(#3\right)}}
% Rectangular full swap walk
% Args:
%       #1='k' size of initial face s
%       #2='l' size of final face t
\newcommand{\rfswap}[2]{\Ess_{#1, #2}}
% Height Specific Squared Swap Walk
% Args:
%       #1='k' size of initial face s
%       #2='j' number of new elements, i.e., |t\s|
%       #3='u' walk height
\newcommand{\hswap}[3]{\hrswap{#1}{#1}{#2}{#3}}

% Squared Swap Walk (Height Independent)
% Args:
%       #1='k' size of initial face s
%       #2='j' number of new elements, i.e., |t\s|
\newcommand{\swap}[2]{\rswap{#1}{#1}{#2}}


% Walk Swap Operator
% Args:
%      #1='k_1'
%      #2='k_2'
\newcommand{\wswap}[2]{\Ess_{#1,#2}^{\circ}}

\newcommand{\tswap}[3]{\Ess_{#1,#2,#3}}

\newcommand{\fswap}[3]{\mathfrak{S}_{#1,#2,#3}}

\newcommand{\sswap}[1]{\Ess_{#1,#1}^{\bigtriangleup}}

%%%%%%%%%%%%%%%%%%%%%%%%%%%%%%%%
%        Normal Walks
%%%%%%%%%%%%%%%%%%%%%%%%%%%%%%%%

% Rectangular Normal Walk
% Args:
%       #1='k' size of initial face
%       #2='l' size of final face
%       #3='u' walk height
\newcommand{\rnwalk}[3]{\Enn_{#1,#2}^{\left(#3\right)}}

% Normal Walk
% Args:
%       #1='k' size of initial face
%       #3='u' walk height
\newcommand{\nwalk}[2]{\Enn_{#1,#1}^{\left(#2\right)}}

\newcommand{\lict}[1]{$(1/2-#1)$-\text{close}}

%%%%%%%%%%%%%%%%%%%%%%%%%%%%%%%%%%%%%%%%%%%%%%%%%%%%%%%%%%%%%%%%%%%%%%%%
% Commands from
% pages.cs.wisc.edu/~dieter/Courses/2013s-CS880/Scribes/TEX/lecture.sty
%%%%%%%%%%%%%%%%%%%%%%%%%%%%%%%%%%%%%%%%%%%%%%%%%%%%%%%%%%%%%%%%%%%%%%%%
% circle a character
\usepackage{tikz}
\newcommand*\circled[1]{\tikz[baseline=(char.base)]{
            \node[shape=circle,draw,inner sep=1pt] (char) {#1};}}
% replacement product
\DeclareMathOperator{\replacement}{\circled{{\rm r}}}
% zig-zag product
\DeclareMathOperator{\zigzag}{\circled{{\rm z}}}
% derandomized squaring
\DeclareMathOperator{\dersquare}{\circled{{\rm s}}}

\newcommand{\gapLarge}{97/100}
\newcommand{\gapSmall}{2/9}
\newcommand{\gapNaive}{7/9+e^{-\poly(n)}} % 1 - gapSmall


\newcommand{\SEP}{\textup{SEP}}
\newcommand{\papo}{\textup{PAPO}}

\newcommand{\Symm}[2]{\vee^{#1}(\C^{#2})}

\usepackage{tikz,float,graphicx}
\usetikzlibrary{shapes.symbols,positioning}
\usepackage[T1]{fontenc}
\setcounter{secnumdepth}{2}
\setcounter{tocdepth}{2}
\usepackage{calc}
\usepackage{amsmath}
\usepackage{amsthm}
\usepackage{amssymb}
\PassOptionsToPackage{normalem}{ulem}
\usepackage{ulem}
\usepackage[font=scriptsize]{caption}

\setenumerate{label=\rm (\roman*)}
\setlist[enumerate,1]{itemsep=-1\parsep, topsep=0mm}
\setlist[description]{itemsep=-1\parsep, topsep=0mm}
\setlist[itemize]{itemsep=-1\parsep, topsep=0mm}

\newcommand{\I}{\mathbf{I}}
\DeclareMathOperator*{\Exp}{\mathbb{E}}
\newcommand{\trace}{\operatorname{Tr}}
\newcommand{\bra}[1]{\langle#1|}
 

\newcommand{\posR}{\mathbb R_{\ge0}}
\newcommand{\TD}{\mathrm{D}}
\newcommand{\SSS}{\mathcal{S}}
\newcommand{\fI}{\mathfrak{I}}
\newcommand{\unitC}{\mathbb{S}}
\newcommand{\unitR}{\mathbb{S}^+}   
\renewcommand{\FF}{\mathbb{F}}
\newcommand{\HH}{\mathbb{H}}
\newcommand{\KK}{\mathbb{K}}
\newcommand{\ZZ}{\mathbb{Z}}

\newcommand{\cM}{\mathcal{M}}
\newcommand{\cL}{\mathcal{L}}
\newcommand{\cR}{\mathcal{R}}
\newcommand{\cV}{\mathcal{V}}
\newcommand{\cK}{\mathcal{K}}

\newcommand{\pred}{\mathcal C}
\newcommand{\cons}{\mathfrak C}


\newcommand{\qpsi}[1]{|\psi_{#1}\rangle}

\newcommand{\qPsi}{|\Psi\rangle}
\newcommand{\qphi}{|\phi\rangle}
\newcommand{\qPhi}{|\Phi\rangle}


%complexity class
%\newcommand{\NEXP}{\mathbf{NEXP}}
%\newcommand{\PCP}{\mathbf{PCP}}
%\newcommand{\Pclass}{\mathbf{P}}



\numberwithin{equation}{section}


\newcommand{\size}{\mathrm{size}}
\newcommand{\succSAT}{\mathbf{succSAT}}
\renewcommand{\outer}{\mathrm{outer}}
\newcommand{\inner}{\mathrm{inner}}

\makeatletter
\newcommand*{\rom}[1]{\expandafter\@slowromancap\romannumeral #1@}
\newcommand{\pushright}[1]{\ifmeasuring@#1\else\omit\hfill$\displaystyle#1$\fi\ignorespaces}
\newcommand{\specialcell}[1]{\ifmeasuring@#1\else\omit$\displaystyle#1$\ignorespaces\fi}
\makeatother


%\newcounter{algcnt}[section]
%\counterwithin*{algcnt}{section}
\newenvironment{myalg}[1][]
{\begin{minipage}[t]{1\columnwidth - 2\fboxsep - 2\fboxrule}
\refstepcounter{theorem}%
\uline{\textbf{Algorithm \thetheorem}: \ignorespaces#1}  \vspace{1mm}\\
}
{\end{minipage}
\vspace{3mm}
}


\newcommand{\pt}{\mathrm{P}}
%\newcommand{\poly}{\mathrm{poly}}

\newcommand{\pei}[1]{{\color{red}#1}}
\newcommand{\sep}{\mathrm{SEP}}
\renewcommand{\d}{\mathrm{d}}

\usepackage{subcaption}
\newcommand{\upth}{\textsuperscript{th}}


\LinesNumberedHidden

\newcommand{\fnote}[1]{\textcolor{blue}{{\textbf{(Fernando: #1)}}}}
\newcommand{\pnote}[1]{\textcolor{olive}{{\textbf{Pei: #1}}}}
